\chapter{Conclusion}\label{ch:conclusion}

\section{The arc, restated}\label{sec:conc-arc}

This thesis began from an observation about changing economics: as capable
code-generating agents drive the marginal cost of a plausible implementation
towards zero, the durable, scarce artefact of a software project ceases to be the
code and becomes the precise, machine-checkable statement of what the code must
do. Formal specifications have always promised to be that statement, and the
machinery for \emph{checking} programs against them---design by contract
\citep{meyer1992dbc}, the Java Modeling Language \citep{leavens2006jml}, and
verifiers such as OpenJML \citep{cok2014openjml} discharging obligations through
Z3 \citep{demoura2008z3}---has been mature for years. What was missing was an
economical supply of the contracts themselves, and a place for them to do useful
work across the lifecycle.

The programme answered that gap in three movements, and its shape is captured by
its title. \emph{Inference} makes specifications cheap: a heuristic static engine
recovers sound, verifiable contracts from existing code at negligible authoring
cost. \emph{Composition} makes them reach: a weakest-precondition pass propagates
those contracts across a call graph, and the same propagation, run across two
versions of a dependency, turns behavioural version compatibility into a static,
run-free analysis. \emph{Integration} makes them load-bearing: the inferred,
verified, compositional contract becomes the verification spine of an AI-native
development lifecycle---at once the input an agent builds against and the oracle
its output is checked against. Inference manufactures the artefact, composition
extends its scope, and integration gives it somewhere to earn its keep; none of
the three stands alone.

\section{Contributions and research questions}\label{sec:conc-contributions}

The three contributions map onto the three research questions answered in
\cref{ch:discussion}.

The first contribution (\cref{ch:inferrer}) is the JML-Inferrer: a static engine
that recovers JML preconditions, postconditions, loop invariants, and assignable
clauses from unannotated Java by pattern-matching on the abstract syntax tree.
Its output is validated end to end against OpenJML, so an emitted contract is not
merely plausible but demonstrably discharges under extended static checking. A
controlled comparison then shows that supplying these contracts to a language
model measurably improves the unit tests it generates, relative to the same model
working from code alone---because a contract supplies the independent oracle a
generated test needs \citep{barr2015oracle}. This answers RQ1 affirmatively: even
heuristic, partial, automatically inferred specifications are good enough to
deliver downstream utility.

The second contribution (\cref{ch:compositional}) is a compositional
weakest-precondition pass that strictly extends single-pass inference. Propagating
a callee's contract into its callers recovers well-formed preconditions the local
pass cannot---approximately $18{,}854$ additional preconditions on the principal
corpus, reported after well-formedness filtering, so vacuous or unsatisfiable-scope
contracts are excluded from the count. Because the reasoning that recovers a
caller's obligation from a callee's contract is exactly the reasoning that
propagates a contract \emph{change}, the identical mechanism, run differentially,
computes the behavioural blast-radius of a dependency upgrade. This answers RQ2: a
dependency change becomes a contract-diff event whose effect on clients can be
checked rather than discovered in production---a behavioural notion of
compatibility finer than the declared semantics of a version number
\citep{preston2013semver}.

The third contribution (\cref{ch:integration}) is a verification-centric
reference architecture in which one machine-checkable contract plays a dual role:
input to the implementing agent and oracle for the verifier, with the verifier
returning counterexamples that close the loop. It supplies the three things that
current specification-driven agentic lifecycles---the AWS AI-Driven Development
Lifecycle \citep{awsaidlc}, GitHub's Spec Kit \citep{githubspeckit}, and Kiro
\citep{awskiro}---lack: a machine-checkable contract rather than prose intent, a
closed counterexample-guided loop, and an oracle whose authority is independent of
the code-writing agent. This answers RQ3 at the level of a design validated by
instantiation of its components.

\section{Limitations}\label{sec:conc-limitations}

Four limitations bound these claims, and the thesis is deliberate about each. The
first is the boundary between a \emph{verified} specification, which the checker
has discharged, and a merely \emph{well-formed} one, which is syntactically valid
and scope-correct but has not passed that check; the compositional reach figure is
a count of well-formed preconditions, not of verified ones, and the two are never
conflated. The second is the partiality of heuristic inference: the engine
recognises the structural patterns it was built for and stays silent elsewhere, so
its fidelity is a rate over emitted contracts, not a coverage of all behaviour.
The third is the reach of the verification substrate: Z3 timeouts, non-linear
arithmetic, and quantified array theories place some true specifications beyond
mechanical discharge, so silence from the checker is never a guarantee of safety.
The fourth is that the reference pipeline of \cref{ch:integration} is designed and
its components individually evidenced, but the integrated, agent-in-the-loop system
has not been subjected to a controlled end-to-end trial. A fifth caveat cuts
across the empirical results: the strongest quantitative findings rest on a single
principal corpus, which establishes the effects in a realistic setting without
establishing their generality.

\section{Future work}\label{sec:conc-future}

Three directions follow directly. The first is industrial validation of the
AI-native lifecycle: an action-research case study of the specification-centred
workflow in a genuine development team \citep{runeson2009casestudy,yin2018case},
framed to claim feasibility and realism rather than a controlled effect size,
subject to Human Research Ethics Committee approval and to careful handling of
confounds such as team experience and the simultaneous adoption of other tooling.
The second is extending the compositional analysis---deeper inter-procedural
reasoning and tighter modelling of standard-library and framework contracts---to
raise the fraction of genuine incompatibilities the gate detects, moving it from
sound-by-construction towards sound-and-comprehensive, and to replicate the
findings on codebases beyond the principal corpus, including the messy legacy
systems that most need retrofitted specifications \citep{feathers2004legacy}. The
third is the direction deliberately deferred throughout: LLM-assisted inference at
inference time, in which a language model proposes candidate contracts that the
formal toolchain then certifies or rejects. The thesis has held this at arm's
length precisely so that its fidelity claims rest on deterministic heuristics and
mechanical verification; introducing a generative proposer while preserving that
guarantee---the verifier remains the authority, the model only a source of
candidates---is a natural and promising extension.

\section{Closing}\label{sec:conc-closing}

When an agent can regenerate an implementation cheaply, the implementation stops
holding the value; what endures is the statement of what the software must do. This
thesis has shown how to make that statement cheap, through inference that recovers
verified contracts from existing code; how to make it reach, through compositional
propagation that reasons about compatibility across a call graph; and how to make
it load-bearing, as the verification spine of how software is built. The path from
inference to integration is, in the end, a single argument: a formal specification,
made cheap and placed at the centre, is the right foundation for software
development in an age when the code can write itself but the intent still cannot.
